\documentclass[11pt]{article}
\usepackage[utf8]{inputenc}
\usepackage[letterpaper, margin=0.6in]{geometry}
\usepackage{amsmath}
\usepackage{amssymb}
\usepackage{amsfonts}
\usepackage{enumitem}
\usepackage{chemfig}
\usepackage{tabularx}

% Custom command for the top header box
\newcommand{\headerbox}{
    \noindent\begin{tabular}{|p{11cm}|p{5.5cm}|}
    \hline
    \scriptsize \textbf{Communication:} & \scriptsize \textbf{Overall:}\\
    \scriptsize Language: easy to understand, answers question, follows instructions, concise, precise & \\
    \scriptsize Chemistry: proper vocab, notation, spacing, logical & \scriptsize \space\space \space\space vague/basic/OK/good/clearly written \\
    \scriptsize Math: organized, units, sig digs, math notation, concluding statement, spacing & \\
    \hline
    \end{tabular}
    \vspace{1em}
}

\begin{document}

% --- PAGE 1 ---
\headerbox

\noindent SCH4U \hfill Name: \underline{\hspace{5cm}} \\
\indent \hfill Monday September 29, 2025 \vspace{1em}

\begin{center}
    \Large \textbf{Unit 1: Structures \& Properties Test (AM)}
\end{center}

\begin{center}
    \begin{tabular}{|lc|lc|lc|}
    \hline
    LL: & \space\space /4 & Comm: & \space\space /10 & Total: & \space\space /40 \\
    \hline
    \end{tabular}
\end{center}
\vspace{1em}

\begin{enumerate}[leftmargin=*, label=\arabic*.]
    \item {[2 marks]} Agree or disagree with this statement \& explain why:
    \begin{center}
        \textit{"Hydrogen only produces 4 lines in its emission spectrum"}
    \end{center}
    \vspace{6em}

    \item {[3 marks]} Which of the following sets of quantum numbers ($n, \ell, m_\ell, m_s$) are possible? If it's possible, which element's last electron has this quantum number?
    \begin{enumerate}[label=\textbf{\alph*.)}, itemsep=2em, topsep=1em]
        \item $4, 0, 0, +\frac{1}{2}$ \hspace{3cm} \textbf{b.)} $3, 2, -3, +\frac{1}{2}$ \hspace{3cm} \textbf{c.)} $2, 3, 0, -\frac{1}{2}$
    \end{enumerate}
    \vspace{4em}

    \item {[2 marks]} Suggest a possible shorthand electron configuration for neutral Rh (element 45). Is this configuration anomalous? Explain.
    \vspace{6em}

    \item {[3 marks]} Using formal charge, identify which structure is \textbf{least} likely. Explain your choice.
    \vspace{1em}
    
    \indent \chemfig{\charge{180=\:}{N}~N-\charge{90=\:,0=\:,270=\:}{O}} \vspace{1.5em}
    
    \indent \chemfig{\charge{180=\:,90=\:}{N}=N=\charge{90=\:,0=\:}{O}} \vspace{1.5em}
    
    \indent \chemfig{\charge{180=\:,90=\:,270=\:}{N}-N~\charge{90=\:}{O}}
    \vspace{6em}
\end{enumerate}

\pagebreak

% --- PAGE 2 ---
\headerbox

\begin{enumerate}[leftmargin=*, label=\arabic*., start=5]
    \item For the element molybdenum (element 42):
    \begin{enumerate}[label=\textbf{\alph*.}]
        \item {[2 marks]} Predict the most stable cation for molybdenum. Explain briefly.
        \vspace{5em}
        \item {[2 marks]} Draw the energy level diagram of the most stable cation from (a). Label all subshells.
        \vspace{1.5em}
        \begin{center}
            \begin{tabular}{l l l l}
                \rule{1.5cm}{0.4pt} & & & \\
                & & & \\
                \rule{1.5cm}{0.4pt} & \rule{0.6cm}{0.4pt} \space \rule{0.6cm}{0.4pt} \space \rule{0.6cm}{0.4pt} & & \\
                & & & \\
                \rule{1.5cm}{0.4pt} & \rule{0.6cm}{0.4pt} \space \rule{0.6cm}{0.4pt} \space \rule{0.6cm}{0.4pt} & \rule{0.5cm}{0.4pt} \space \rule{0.5cm}{0.4pt} \space \rule{0.5cm}{0.4pt} \space \rule{0.5cm}{0.4pt} \space \rule{0.5cm}{0.4pt} & \\
                & & & \\
                \rule{1.5cm}{0.4pt} & \rule{0.6cm}{0.4pt} \space \rule{0.6cm}{0.4pt} \space \rule{0.6cm}{0.4pt} & \rule{0.5cm}{0.4pt} \space \rule{0.5cm}{0.4pt} \space \rule{0.5cm}{0.4pt} \space \rule{0.5cm}{0.4pt} \space \rule{0.5cm}{0.4pt} & \rule{0.5cm}{0.4pt} \space \rule{0.5cm}{0.4pt} \space \rule{0.5cm}{0.4pt} \space \rule{0.5cm}{0.4pt} \space \rule{0.5cm}{0.4pt} \\
                & & & \\
                \rule{1.5cm}{0.4pt} & & & \\
            \end{tabular}
        \end{center}
    \end{enumerate}
    \vspace{5em}

    \item Some atoms have anomalous electron configurations:
    \begin{enumerate}[label=\textbf{\alph*.}]
        \item {[2 marks]} List all the atoms with anomalous configurations in the 3d subshell.
        \vspace{4em}
        \item {[2 marks]} Is silicon likely to have an anomalous configuration? Explain.
        \vspace{6em}
    \end{enumerate}

    \item {[2 marks]} Why is anomalous electron configuration not the same as hybridization? Provide an example.
    \vspace{8em}
\end{enumerate}

\pagebreak

% --- PAGE 3 ---
\headerbox

\begin{enumerate}[leftmargin=*, label=\arabic*., start=8]
    \item {[4 marks]} Draw the orbital overlap showing the bonding in the following structure. Use different colours for each bond (if you have them). Label each bond as being $\sigma$ or $\pi$.
    \vspace{2em}
    
    \indent \chemfig{H-Si(-[6]F)(=[1]O)}
    \vspace{12em}

    \item {[2 marks]} What is the same \& what is different between the following: $AB_2E$ \& $AB_2E_2$? Be precise.
    \vspace{10em}

    \item {[2 marks]} Why are ionic solids brittle? What is unique about the structure that allows it to be so? Provide an example of an ionic solid.
    \vspace{10em}

    \item For the compound $SF_4$
    \begin{enumerate}[label=\textbf{\alph*.}]
        \item {[1 mark]} Write the VSEPR notation \hfill \underline{\hspace{5cm}}
        \item {[1 mark]} Draw the 3D ball \& stick model
        \vspace{10em}
        \item {[2 marks]} Is the molecule polar? Explain
        \vspace{6em}
        \item {[1 mark]} Circle the intermolecular force(s) \hfill LF \space \space \space \space dd \space \space \space \space H-bond
    \end{enumerate}
\end{enumerate}

\pagebreak

% --- PAGE 4 ---
\headerbox

\begin{enumerate}[leftmargin=*, label=\arabic*., start=12]
    \item {[2 marks]} Element X releases red light when an e falls from $n = 5$ to $n = 3$. Element Y releases violet light when an electron falls from $n = 5$ to $n = 3$. Which element would you predict has the larger radius from nucleus to the outermost electrons? Explain briefly.
    \vspace{12em}
\end{enumerate}

\noindent \underline{\textbf{Multiple Choice:}} \\
Put your answers to the multiple choice in the table. \\
\noindent \begin{tabularx}{\textwidth}{|X|X|X|X|X|}
\hline
1. & 2. & 3. & 4. & 5. \\
\hline
\rule{0pt}{2em} & & & & \\
\hline
\end{tabularx}
\vspace{2em}

\begin{enumerate}[leftmargin=*, label=\arabic*.]
    \item Which of the following represents the ground state electron configuration of the $Zn^{2+}$ ion?
    \begin{enumerate}[label=\textbf{\Alph*.}]
        \item $1s^2 2s^2 2p^6 3s^2 3p^6 4s^2 3d^{10} 4p^2$ \hfill \textbf{D.} $1s^2 2s^2 2p^6 3s^2 3p^6 3d^8 4p^2$
        \item $1s^2 2s^2 2p^6 3s^2 3p^6 3d^{10}$ \hfill \textbf{E.} $1s^2 2s^2 2p^6 3s^2 3p^6 4s^2 3d^8$
        \item $1s^2 2s^2 2p^6 3s^2 3p^6 4s^2 4p^6 3d^8$
    \end{enumerate}
    \vspace{1em}

    \item Which statement about bonding orbitals is FALSE?
    \begin{enumerate}[label=\textbf{\Alph*.}]
        \item the overlap of two s-orbitals forms a $\sigma$-bond
        \item the overlap of one $sp^2$ orbital \& one p-orbital forms a $\pi$-bond
        \item the overlap of two $sp^2$ orbitals forms a $\sigma$-bond
        \item the overlap of two p-orbitals form a $\sigma$-bond
        \item all are false
    \end{enumerate}
    \vspace{1em}

    \item In a molecule of $C_2H_4$, what types of orbitals form bonds around each \textbf{carbon} atom?
    \begin{enumerate}[label=\textbf{\Alph*.}]
        \item two s-orbitals \& two p-orbitals \hfill \textbf{D.} one p-orbital \& three $sp^3$-orbitals
        \item four $sp^3$ orbitals \hfill \textbf{E.} one p-orbital \& three $sp^2$-orbitals
        \item four $sp^2$ orbitals
    \end{enumerate}
    \vspace{1em}

    \item Which of the following molecules have polar bonds but are non-polar molecules? \\
    
    \noindent \textbf{I.} $BeH_2$ \hspace{1cm} \textbf{II.} $BH_3$ \hspace{1cm} \textbf{III.} $H_2O$ \hspace{1cm} \textbf{IV.} $NH_3$ \hspace{1cm} \textbf{V.} $CCl_4$ \vspace{0.5em}
    
    \noindent
    \begin{tabularx}{\textwidth}{XXXXX}
        \textbf{A.} I \& II & \textbf{B.} I \& V & \textbf{C.} II \& III & \textbf{D.} II, III \& IV & \textbf{E.} I, II \& V
    \end{tabularx}
    \vspace{1em}

    \item What is the correct notation for a molecule containing a central atom attached to three other atoms with two lone pairs of electrons?
    
    \begin{tabularx}{\textwidth}{XXXXX}
        \textbf{A.} $AB_5$ & \textbf{B.} $AB_3E_2$ & \textbf{C.} $AB_2E_3$ & \textbf{D.} $AE_5$ & \textbf{E.} $AB_5E_3$
    \end{tabularx}
\end{enumerate}

\end{document}